% p379_1; source PDF page 379; manually verified editable table.
\begin{tikzpicture}[x=1cm,y=-1cm,inner sep=0pt,outer sep=0pt]
\path[use as bounding box] (0,0) rectangle (16.000,3.745);
\draw[line width=.4pt] (0.000,0.000) rectangle (4.950,1.115);
\node[align=center,text width=4.690cm,font={\linespread{1}\fontsize{10}{13.50}\selectfont\bfseries}] at (2.475,0.557) {案件模型};
\draw[line width=.4pt] (4.950,0.000) rectangle (9.900,1.115);
\node[align=center,text width=4.690cm,font={\linespread{1}\fontsize{10}{13.50}\selectfont\bfseries}] at (7.425,0.557) {正常销赃，未欺骗买家};
\draw[line width=.4pt] (9.900,0.000) rectangle (16.000,1.115);
\node[align=center,text width=5.840cm,font={\linespread{1}\fontsize{10}{13.50}\selectfont\bfseries}] at (12.950,0.557) {欺骗买家，谎称是自己的财物};
\draw[line width=.4pt] (0.000,1.115) rectangle (4.950,2.180);
\node[align=center,text width=4.690cm,font={\linespread{1}\fontsize{10}{13.50}\selectfont}] at (2.475,1.647) {甲偷了丙的财物，甲销赃};
\draw[line width=.4pt] (4.950,1.115) rectangle (9.900,2.180);
\node[align=left,text width=4.690cm,font={\linespread{1}\fontsize{10}{13.50}\selectfont}] at (7.425,1.647) {不可罚的事后行为};
\draw[line width=.4pt] (9.900,1.115) rectangle (16.000,2.180);
\node[align=left,text width=5.840cm,font={\linespread{1}\fontsize{10}{13.50}\selectfont}] at (12.950,1.647) {定诈骗罪，与盗窃罪并罚};
\draw[line width=.4pt] (0.000,2.180) rectangle (4.950,3.745);
\node[align=center,text width=4.690cm,font={\linespread{1}\fontsize{10}{13.50}\selectfont}] at (2.475,2.962) {甲偷了丙的财物，乙帮甲销赃};
\draw[line width=.4pt] (4.950,2.180) rectangle (9.900,3.745);
\node[align=left,text width=4.690cm,font={\linespread{1}\fontsize{10}{13.50}\selectfont}] at (7.425,2.962) {定掩饰、隐瞒犯罪所得罪};
\draw[line width=.4pt] (9.900,2.180) rectangle (16.000,3.745);
\node[align=left,text width=5.840cm,font={\linespread{1}\fontsize{10}{13.50}\selectfont}] at (12.950,2.962) {同时触犯掩饰、隐瞒犯罪所得罪和诈骗罪，想象竞合，择一重罪论处};
\end{tikzpicture}
